\uuid{Lx5M}
\titre{Congruences}
\auteur{bodin ; croizat ; sacre}
\organisation{exo7}
\licence{CC BY-NC-SA 4.0}
\datecreate{2022-03-08}
\niveau{L1}
\module{Algèbre}
\chapitre{Arithmétique dans Z}
\sousChapitre{Autre}
\difficulte{3}
\theme{Congruences}
\exo7qcmid{hilisit.arithmetique.050}

\contenu{
\begin{question}
\qclassification{103.99}
 Soit un entier $k$ tel que $k \equiv 2 \;[7]$. Quelles sont les affirmations vraies ?
    \begin{answers}
        \bad{$2k^2 + k \equiv k^3 \;[7]$}
        \good{$3(k^4-k) \equiv 0 \;[7]$}
        \good{$14k-2 \equiv 5 \;[7]$}
        \bad{$k^{18} + k^{12} + k^6 \equiv k \;[7]$}
    \end{answers}
    \begin{explanations}
    On calcule que $k^2 \equiv 4 \;[7]$ ; $k^3 \equiv 2^3 \equiv 1 \;[7]$ et
    $k^4 \equiv 2^4 \equiv 2 \;[7]$. On a alors :\\
    $2k^2 + k \equiv 2 \times 2^2 + 2 \equiv 10 \equiv 3 \;[7]$.\\
    $k^4 \equiv k \;[7]$ donc $3(k^4-k) \equiv 3 \times 0 \equiv 0 \;[7]$.\\
    $14k \equiv 7 \times 2k \equiv 0 \;[7]$ donc $14k-2 \equiv -2 \equiv 5 \;[7]$.\\
    Enfin on a $k^{18} \equiv k^{12} \equiv k^6 \equiv 1 \;[7]$ (c'est le théorème de Fermat, ou une conséquence directe de $k^3 \equiv 1 \;[7]$). Donc
    $k^{18} + k^{12} + k^6 \equiv 3 \;[7]$.
    \end{explanations}
\end{question}
}
